\unitopen{35}
\begin{candos}
\item nazwać główne instytucje i dokumenty, z którymi zetkniesz się po przeprowadzce do Szwajcarii;
\item zrozumieć różnice między niemieckim standardowym a szwajcarskim niemieckim w słownictwie i pisowni;
\item załatwić zameldowanie, ubezpieczenie i telefon, znając potrzebne słowa;
\item rozpoznać dwadzieścia najczęstszych słów dialektalnych.
\end{candos}

\section{Słownictwo}
\vocabtable{35}

\begin{uwaga}[Ta jednostka jest wyjątkowo wrażliwa na aktualność]
Prawo, opłaty i terminy w Szwajcarii zmieniają się i zależą od kantonu i gminy. Wszystko, co oznaczono \verify{tak}, to
twierdzenie, które powinieneś potwierdzić w urzędowym źródle (ch.ch, strona Twojej gminy i kantonu, SEM) przed decyzją.
\end{uwaga}

\section{Pierwszy miesiąc w Szwajcarii}
\begin{center}\small
\begin{tabular}{@{}p{0.25\linewidth}p{0.67\linewidth}@{}}
\toprule
\textsf{\textbf{Sprawa}} & \textsf{\textbf{Co robisz}}\\
\midrule
Zameldowanie & zgłaszasz się do \de{Einwohnerkontrolle} swojej gminy w krótkim terminie od przyjazdu \verify{(zwykle ok.\ 14 dni; sprawdź w gminie)}\\
Pozwolenie & obywatel UE/EFTA z umową o pracę otrzymuje zwykle \de{Aufenthaltsbewilligung B} \verify{(sprawdź warunki w SEM)}\\
Ubezpieczenie zdrowotne & obowiązkowe; wybierasz \de{Krankenkasse}, \de{Franchise} i \de{Prämie} \verify{(termin na wybór: około trzech miesięcy)}\\
Konto i telefon & \de{Bankkonto}, \de{Handyabo} (CH: \de{Mobile-Abo})\\
Podatki & pracownik zagraniczny często płaci \de{Quellensteuer} \verify{(zależy od pozwolenia i dochodu)}\\
Emerytura & \de{AHV} (1.\ filar), \de{Pensionskasse} (2.\ filar), opcjonalnie \de{3.\ Säule} \verify{(sprawdź zasady dla zagranicznych pracowników)}\\
Odpady & worki opłacane gminnie (\de{Kehrichtsack}), segregacja jest obowiązkowa; cisza nocna zwykle od 22 \verify{(zależy od gminy)}\\
\bottomrule
\end{tabular}
\end{center}

\section{Standard i dialekt: różnice, które widać od razu}
\begin{center}\small
\begin{tabular}{@{}lll@{}}
\toprule
\textsf{\textbf{Standard (DE)}} & \textsf{\textbf{Standard CH}} & \textsf{\textbf{Dialekt (Zürich)}}\\
\midrule
\de{Fahrrad} & \de{Velo} & \de{Velo}\\
\de{Fahrkarte} & \de{Billett} & \de{Billett}\\
\de{parken} & \de{parkieren} & \de{parkiere}\\
\de{Bürgersteig} & \de{Trottoir} & \de{Trottoir}\\
\de{Guten Tag} & \de{Grüezi} & \de{Grüezi}\\
\de{Danke} & \de{Merci} & \de{Merci}\\
\de{Tschüss} & \de{Ade, Adieu} & \de{Ade, Uf Wiederluege}\\
\de{Mittagessen} & \de{Mittagessen} & \de{Zmittag}\\
\de{Frühstück} & \de{Frühstück} & \de{Zmorge}\\
\de{Straße} & \de{Strasse} & \de{Strass}\\
\bottomrule
\end{tabular}
\end{center}
Pisownia: \emph{Szwajcaria nie używa \de{ß}}: zawsze \de{ss}. Dialekt jest mową, nie pismem: pisze się po standardowemu, mówi po dialektalnemu. Końcówka
zdrobniała \de{-li} (\de{Brötli, Müesli}) odpowiada \de{-chen}. \de{Gäll}/\de{oder?} na końcu to „prawda?”.

\begin{falle}
Typowy błąd: ucząc się dialektu \emph{przed} B1, odkładasz standard. Szwajcarzy przechodzą na standard, kiedy mówią do Ciebie jak do
uczącego się. Cel: \emph{rozumieć} dialekt (słuchaj SRF, \de{Schawinski}, podcasty), mówić standardem, i nie przepraszać za to. Druga pułapka:
\de{Grüezi} jest uprzejme także do jednej osoby; do wielu \de{Grüezi mitenand}.
\end{falle}

\begin{szwajcaria}
Szwajcaria ma cztery języki narodowe (niemiecki, francuski, włoski, romansz), więc w Zurychu urzędy i firmy używają niemieckiego, a w
Genewie francuskiego. Szwajcarzy cenią prywatność, porządek i ciszę (\de{Ruhezeit}), a spotkanie towarzyskie zaczyna się często od
\de{Apéro}. Zaproszenie na \de{Apéro} jest okazją do poznania ludzi, nie obowiązkiem jedzenia.
\end{szwajcaria}

\begin{praca}
Sprawy kadrowe: \de{Ich habe meine Aufenthaltsbewilligung erhalten.} \de{Ich muss mich noch bei der Einwohnerkontrolle anmelden.}
\de{Wer ist für die Pensionskasse zuständig?} \de{Wie viel Franchise soll ich wählen?}
\end{praca}

\section{Dialogi}
\textbf{Dialog 1 --- w urzędzie gminy}
\begin{dialog}
\turn{A}{Grüezi. Ich möchte mich anmelden. Ich bin gestern nach Zürich gezogen.}{Dzień dobry. Chcę się zameldować. Wczoraj przeprowadziłem się do Zurychu.}
\turn{B}{Haben Sie Ihren Arbeitsvertrag und Ihren Pass dabei?}{Czy ma Pan umowę o pracę i paszport?}
\turn{A}{Ja, hier sind sie. Brauche ich auch den Mietvertrag?}{Tak, oto one. Czy potrzebuję też umowy najmu?}
\turn{B}{Ja, bitte. Dann kann ich Sie anmelden.}{Tak, proszę. Wtedy mogę Pana zameldować.}
\end{dialog}

\section{Ćwiczenia}
\begin{cwiczenia}
\exitem{Podaj szwajcarski odpowiednik: \de{Fahrrad}.}{\de{Velo}}
\exitem{Podaj szwajcarski odpowiednik: \de{Fahrkarte}.}{\de{Billett}}
\exitem{Podaj szwajcarski odpowiednik: \de{parken}.}{\de{parkieren}}
\exitem{Napisz po szwajcarsku (bez \de{ß}): \de{Straße}.}{\de{Strasse}}
\exitem{Jak pozdrowisz kilka osób po szwajcarsku?}{\de{Grüezi mitenand!}}
\exitem{Co znaczy \de{Znüni} i \de{Zvieri}?}{przekąska ok.\ 9:00 i przekąska po południu}
\exitem{Powiedz: „Chcę się zameldować.”}{\de{Ich möchte mich anmelden.}}
\exitem{Zapytaj: „Czy mogę obejrzeć mieszkanie w czwartek?”}{\de{Kann ich die Wohnung am Donnerstag besichtigen?}}
\exitem{Jaki jest cel \de{Franchise}?}{wysokość rocznego udziału własnego w kosztach leczenia (im wyższy, tym niższa składka)}
\exitem{Popraw pisownię CH: \de{Wir haben Straße gelernt}.}{\de{Wir haben Strasse gelernt} (CH: bez \de{ß})}
\exitem{Przetłumacz: „Rozumiem dialekt, ale mówię po standardowemu.”}{\de{Ich verstehe den Dialekt, aber ich spreche Hochdeutsch.}}
\end{cwiczenia}

\weekend{Sporządź listę kontrolną przeprowadzki do Szwajcarii (dokumenty, urzędy, ubezpieczenia) po niemiecku w dziesięciu punktach,
każdy z zadaniem i terminem. Sprawdź każdą datę i opłatę w źródle urzędowym, a nie w tej książce. Wysłuchaj dwóch minut SRF i
wypisz pięć słów, które były nowe.}
